\documentclass[11pt,a4paper]{article}

\usepackage[utf8]{inputenc}
\usepackage[T1]{fontenc}
\usepackage[ngerman]{babel}

\usepackage{amsmath,amssymb,amsthm,mathtools}
\usepackage[a4paper,margin=2.4cm]{geometry}
\usepackage{lmodern}
\usepackage{microtype}
\usepackage{hyperref}
\usepackage{enumitem}

\title{Eine straffe lokale Strukturtheorie des EABC-Gitters modulo \(12\)}
\author{Kollaborative Ausarbeitung}
\date{\today}

\newtheorem{satz}{Satz}
\newtheorem{lemma}[satz]{Lemma}
\newtheorem{korollar}[satz]{Korollar}
\theoremstyle{definition}
\newtheorem{definition}[satz]{Definition}
\newtheorem{beispiel}[satz]{Beispiel}
\theoremstyle{remark}
\newtheorem{bemerkung}[satz]{Bemerkung}

\newcommand{\N}{\mathbb{N}}
\newcommand{\pp}{\mathbb{P}}
\newcommand{\ind}{\mathbf{1}}

\begin{document}
	\maketitle
	
	\begin{abstract}
		Wir formulieren eine kompakte lokale Strukturtheorie der kleinen Primzahlkonstellationen
		im modulo-\(12\)-Gitter
		\[
		Q_n=(12n-11,\,12n-7,\,12n-5,\,12n-1).
		\]
		Dabei werden Primzahlen \(>3\), Primzahlzwillinge, Primzahldrillinge,
		blockübergreifende Primzahlvierlinge und dichte Primzahlfünflinge in einem gemeinsamen
		Rahmen beschrieben. Der Hauptsatz zeigt, dass diese Konstellationen im EABC-Gitter
		vollständig durch zwei Zwillingskanäle, einen singulären Drillingsanker und orientierte
		Erweiterungen von Doppelstellen bestimmt werden.
	\end{abstract}
	
	\section{Das EABC-Gitter}
	
	\begin{definition}[EABC-Familien]
		Für \(n\in\N\) definieren wir
		\[
		E_n:=12n-11,\qquad
		A_n:=12n-7,\qquad
		B_n:=12n-5,\qquad
		C_n:=12n-1.
		\]
		Der Viererblock
		\[
		Q_n:=(E_n,A_n,B_n,C_n)
		\]
		heiße der \(n\)-te \emph{EABC-Truppel}.
	\end{definition}
	
	\begin{lemma}[Primzahlen \(>3\) im EABC-Gitter]
		Jede Primzahl \(p>3\) liegt eindeutig in genau einer der vier Familien
		\[
		E,\quad A,\quad B,\quad C.
		\]
	\end{lemma}
	
	\begin{proof}
		Ist \(p>3\) prim, so ist \(p\) weder durch \(2\) noch durch \(3\) teilbar, also
		\[
		p\equiv 1,5,7,\text{ oder }11 \pmod{12}.
		\]
		Diese vier Restklassen werden disjunkt und vollständig durch
		\[
		12n-11,\quad 12n-7,\quad 12n-5,\quad 12n-1
		\]
		parametrisiert.
	\end{proof}
	
	\section{Zwei Zwillingskanäle}
	
	\begin{definition}[Interner und äußerer Kanal]
		Zu \(Q_n\) definieren wir
		\[
		T_n^{\mathrm{int}}:=(A_n,B_n)=(12n-7,\,12n-5),
		\]
		\[
		T_n^{\mathrm{ext}}:=(C_n,E_{n+1})=(12n-1,\,12n+1).
		\]
		Ferner setzen wir
		\[
		\tau(Q_n):=
		\bigl(
		\ind_{T_n^{\mathrm{int}}\text{ ist Primzahlzwilling}},
		\ind_{T_n^{\mathrm{ext}}\text{ ist Primzahlzwilling}}
		\bigr)\in\{0,1\}^2.
		\]
	\end{definition}
	
	\begin{lemma}[Vollständigkeit der beiden Zwillingskanäle]
		Jeder Primzahlzwilling \((p,p+2)\) mit \(p>3\) liegt eindeutig in genau einem der beiden Muster
		\[
		(12n-7,12n-5)
		\qquad\text{oder}\qquad
		(12n-1,12n+1).
		\]
	\end{lemma}
	
	\begin{proof}
		Für \(p>3\) sind modulo \(12\) nur die Restmuster
		\[
		(5,7)\quad\text{oder}\quad(11,1)
		\]
		möglich. Diese entsprechen genau den Paaren
		\[
		(A_n,B_n)
		\quad\text{bzw.}\quad
		(C_n,E_{n+1}).
		\]
		Die Eindeutigkeit des Index folgt aus der linearen Parametrisierung.
	\end{proof}
	
	\section{Drillinge als singuläre Randphänomene}
	
	\begin{definition}[Dichte Primzahldrillinge]
		Unter einem \emph{dichten Primzahldrilling} verstehen wir ein Primzahltripel der Form
		\[
		(p,p+2,p+6)
		\qquad\text{oder}\qquad
		(p,p+4,p+6).
		\]
	\end{definition}
	
	\begin{lemma}[Jeder Drilling enthält die \(3\)]
		Jeder dichte Primzahldrilling enthält notwendig die Primzahl \(3\).
	\end{lemma}
	
	\begin{proof}
		In beiden Mustern ist modulo \(3\) stets mindestens eine der drei Zahlen durch \(3\) teilbar.
		Ist sie prim, so muss sie gleich \(3\) sein.
	\end{proof}
	
	\begin{korollar}[Vollständige Klassifikation der Drillinge]
		Die einzigen dichten Primzahldrillinge sind
		\[
		(3,5,7)
		\qquad\text{und}\qquad
		(3,7,11).
		\]
	\end{korollar}
	
	\begin{proof}
		Nach dem Lemma muss \(3\) enthalten sein. Dann bleiben nur die beiden dichten Realisierungen
		\[
		(3,5,7)\quad\text{und}\quad(3,7,11),
		\]
		und beide sind prim.
	\end{proof}
	
	\section{Doppelstellen und Vierlinge}
	
	\begin{definition}[Doppelstelle]
		Ein EABC-Truppel \(Q_n\) heißt \emph{Doppelstelle}, wenn
		\[
		\tau(Q_n)=(1,1).
		\]
	\end{definition}
	
	\begin{lemma}[Doppelstellen und Vierlinge]
		Ein Truppel \(Q_n\) ist genau dann eine Doppelstelle, wenn
		\[
		(12n-7,\;12n-5,\;12n-1,\;12n+1)
		\]
		ein Primzahlvierling ist.
	\end{lemma}
	
	\begin{proof}
		Die Bedingung \(\tau(Q_n)=(1,1)\) bedeutet, dass sowohl
		\[
		(12n-7,12n-5)
		\]
		als auch
		\[
		(12n-1,12n+1)
		\]
		Primzahlzwillinge sind. Das ist genau die Primheit aller vier Zahlen
		\[
		12n-7,\ 12n-5,\ 12n-1,\ 12n+1,
		\]
		also das Vierlingsmuster
		\[
		p,\ p+2,\ p+6,\ p+8.
		\]
	\end{proof}
	
	\section{Fünflinge als orientierte Erweiterungen}
	
	\begin{definition}[Dichter Primzahlfünfling]
		Ein \emph{dichter Primzahlfünfling} ist ein Primzahlmuster der Form
		\[
		(p,p+2,p+6,p+8,p+12).
		\]
	\end{definition}
	
	\begin{lemma}[Zwei Formen der Fünflinge]
		Jeder dichte Primzahlfünfling mit \(p>3\) besitzt genau eine der beiden Formen
		\[
		(A_n,B_n,C_n,E_{n+1},A_{n+1})
		\]
		oder
		\[
		(C_n,E_{n+1},A_{n+1},B_{n+1},C_{n+1})
		\]
		für ein eindeutig bestimmtes \(n\).
	\end{lemma}
	
	\begin{proof}
		Für \(p>3\) sind als Startreste modulo \(12\) nur \(5\) oder \(11\) möglich.
		
		Falls \(p\equiv 5\pmod{12}\), ist \(p=A_n\), also entsteht
		\[
		(A_n,B_n,C_n,E_{n+1},A_{n+1}).
		\]
		
		Falls \(p\equiv 11\pmod{12}\), ist \(p=C_n\), also entsteht
		\[
		(C_n,E_{n+1},A_{n+1},B_{n+1},C_{n+1}).
		\]
		
		Die Startreste \(1\) und \(7\) sind unmöglich, weil dann eine der Folgezahlen durch \(3\)
		teilbar wäre.
	\end{proof}
	
	\begin{lemma}[Jeder Fünfling enthält genau eine Doppelstelle]
		Jeder dichte Primzahlfünfling enthält genau einen blockübergreifenden Primzahlvierling.
	\end{lemma}
	
	\begin{proof}
		In der ersten Form
		\[
		(A_n,B_n,C_n,E_{n+1},A_{n+1})
		\]
		bilden die ersten vier Zahlen die Doppelstelle.
		
		In der zweiten Form
		\[
		(C_n,E_{n+1},A_{n+1},B_{n+1},C_{n+1})
		\]
		bilden die ersten vier Zahlen ebenfalls die Doppelstelle.
		
		Da ein dichter Fünfling nur fünf Zahlen umfasst, kann er nicht zwei verschiedene Vierlinge
		dieser Form enthalten.
	\end{proof}
	
	\section{Hauptsatz}
	
	\begin{satz}[Lokale Vollständigkeit des EABC-Gitters]
		Das EABC-Gitter
		\[
		Q_n=(12n-11,12n-7,12n-5,12n-1)
		\]
		liefert eine vollständige lokale Beschreibung der niedrigen Primzahlkonstellationen:
		
		\begin{enumerate}[label=\textup{(\roman*)}]
			\item Jede Primzahl \(>3\) liegt eindeutig in einer der vier Familien \(E,A,B,C\).
			\item Jeder Primzahlzwilling \(>3\) liegt eindeutig in genau einem der beiden Kanäle
			\[
			(A_n,B_n)
			\quad\text{oder}\quad
			(C_n,E_{n+1}).
			\]
			\item Jeder dichte Primzahldrilling ist ein singuläres \(3\)-verankertes Randphänomen und gehört
			genau zu
			\[
			(3,5,7)
			\quad\text{oder}\quad
			(3,7,11).
			\]
			\item Jede Doppelstelle \(\tau(Q_n)=(1,1)\) ist genau ein blockübergreifender Primzahlvierling
			\[
			(12n-7,12n-5,12n-1,12n+1).
			\]
			\item Jeder dichte Primzahlfünfling ist genau eine orientierte Erweiterung einer Doppelstelle und
			hat daher genau eine der beiden Formen
			\[
			(A_n,B_n,C_n,E_{n+1},A_{n+1})
			\quad\text{oder}\quad
			(C_n,E_{n+1},A_{n+1},B_{n+1},C_{n+1}).
			\]
		\end{enumerate}
		
		Damit ist die gesamte lokale Zwillings-, Drillings-, Vierlings- und Fünflingsstruktur des
		Primzahlgitters modulo \(12\) vollständig beschrieben.
	\end{satz}
	
	\begin{proof}
		Aussage \textup{(i)} ist Lemma 1. Aussage \textup{(ii)} folgt aus der Vollständigkeit der beiden
		Zwillingskanäle. Aussage \textup{(iii)} ist das Korollar über die vollständige Klassifikation der
		Drillinge. Aussage \textup{(iv)} ist das Lemma über Doppelstellen und Vierlinge. Aussage
		\textup{(v)} folgt aus den beiden Fünflingslemmata.
	\end{proof}
	
	\section{Beispiele}
	
	\begin{beispiel}[Innerer Zwilling]
		Für \(n=2\) ist
		\[
		Q_2=(13,17,19,23),
		\]
		und
		\[
		(A_2,B_2)=(17,19)
		\]
		ist ein Primzahlzwilling.
	\end{beispiel}
	
	\begin{beispiel}[Äußerer Zwilling]
		Für \(n=5\) ist
		\[
		Q_5=(49,53,55,59),
		\]
		und
		\[
		(C_5,E_6)=(59,61)
		\]
		ist ein Primzahlzwilling.
	\end{beispiel}
	
	\begin{beispiel}[Drilling]
		Der Drilling
		\[
		(3,5,7)
		\]
		ist der \(3\)-verankerte Randfall des inneren Kanals.
	\end{beispiel}
	
	\begin{beispiel}[Doppelstelle / Vierling]
		Für \(n=9\) ist
		\[
		Q_9=(97,101,103,107),
		\]
		und
		\[
		(101,103,107,109)
		\]
		ist der zugehörige blockübergreifende Vierling.
	\end{beispiel}
	
	\begin{beispiel}[Fünfling]
		Der Fünfling
		\[
		(5,7,11,13,17)
		\]
		hat die Form
		\[
		(A_1,B_1,C_1,E_2,A_2)
		\]
		und ist die orientierte Erweiterung der Doppelstelle
		\[
		(5,7,11,13).
		\]
	\end{beispiel}
	
	\section{Schluss}
	
	Die modulo-\(12\)-Zerlegung erzeugt im EABC-Gitter eine sehr kompakte lokale Ordnung:
	Zwillingskanäle, singuläre Drillinge, blockübergreifende Vierlinge und orientierte
	Fünflinge erscheinen nicht als isolierte Phänomene, sondern als aufeinander aufbauende
	Konfigurationsstufen desselben Gitters.
	
\end{document}